\section{Cumulative Intellectual Development}
\label{app:development}
The earlier v1 project studied AI attention, interruptions, and reputation-calibrated scheduling. Multiple agents submitted information, but an attention coordinator chiefly scored and scheduled their interruptions. The author concluded that direct strategic interdependence between two decision makers was not sufficiently central or explicit: one agent's optimal action was not clearly contingent on another agent's action, and players, actions, payoffs, best responses, and equilibrium were not the organizing model. This was the author's own reassessment, not a change attributed to either peer review.

The author therefore selected costly information acquisition before a collective AI decision. Two agents directly choose Research or Skip. If the other researches, Skip preserves the shared information benefit without paying the private cost; if the other skips and \(0<c<V\), Research is individually worthwhile. The resulting best responses depend on the opponent's action, making free-riding and equilibrium central. Economics explains the incentives, computation makes the benchmark reproducible through payoff matrices, independent checks, NashPy, a cost sweep, a notebook, tests, and a local prototype, and behavioral science asks whether future observed choices follow that benchmark or differ with beliefs and bounded reasoning. These lenses address one mechanism. The small baseline also permits analytical predictions, computational checks, and explicit separation from future behavioral evidence. Asymmetric costs, signal quality, repeated interaction, larger groups, aggregation, and responsibility mechanisms remain future extensions, not implemented results.

Within the revised game, the author initially predicted the two pure equilibria \((R,S)\) and \((S,R)\). NashPy verified them and additionally identified a symmetric mixed equilibrium. A subsequent indifference calculation yielded \(p(\text{Research})=1-c/V\) for the strict interior \(0<c<V\), agreeing with the tested interior costs. Computation thus revealed equilibrium structure beyond the author's initial pure-strategy reasoning, which was then analytically checked.

The 2056 Nobel and Turing Laureate Program Launch Workshop took place September 7, 2026 at Duke Kunshan University, IB 2050; Prof. Luyao Zhang was Program Chair and Prof. Ken Rogerson was Program Discussant. The author's earlier workshop session was Session B. No workshop comment or peer-review claim is inferred or attributed here. Any course-required handwritten reflection remains the author's own Human-Only work.
